\chapter{Conclusion}\label{cha:conclusion}

\section{Summary of Contributions}\label{sec:summary}

This thesis investigated whether truth-value-preserving
transformations can be used to systematically test QBF solvers by
deliberately activating specific internal mechanisms while
preserving the correct solver output as an oracle. To this end, the
\emph{QBF Transformer} tool was designed and implemented as a
complete experimental pipeline covering parsing, transformation,
solving with DepQBF and CAQE, runtime measurement, and \texttt{gcov}
coverage analysis.

Four operators were formally defined and proven truth-value-preserving
(see Chapter~\ref{cha:operators}): \emph{Tautology Insertion},
\emph{Pure Literal Generator}, \emph{Universal Reduction Trigger},
and \emph{Blocked Clause Insertion}. Each operator targets a
distinct internal solver mechanism.

The empirical evaluation over 285 valid experiments achieved a
100\,\% equivalence rate, confirming that all four operators are
correctly implemented in practice. The aggregated measurements
(Table~\ref{tab:summary}) show that the operators produce
measurable, operator-specific effects on both runtime and coverage,
though the absolute magnitudes are small for the benchmark set
used. A consistent asymmetry between the two solvers emerged:
DepQBF shows runtime ratios at or below~$1.0$ across all operators,
whereas CAQE shows ratios at or above~$1.0$ for most operators.
This suggests solver-specific differences in preprocessing that
warrant further investigation.

\section{Limitations}\label{sec:limitations}

\textbf{Benchmark size.} Only 50 formulas from seven QBFLIB
families were evaluated. The small absolute coverage deltas
(below~$0.10\,\text{pp}$) reflect the limited size of the
benchmarks rather than the absence of an effect; larger and more
diverse sets (e.g.\ the QBF Gallery) would be needed for
statistically robust conclusions.

\textbf{Coverage measurement limited to DepQBF.} \texttt{gcov} is
not directly applicable to CAQE's Rust build, so the runtime
asymmetry observed for CAQE cannot currently be explained at the
code-path level.

\textbf{Equivalence check via DepQBF only.} The equivalence flag
relies on DepQBF's output; a systematic DepQBF bug would not be
detected by this check alone.

\textbf{Structural transformations only.} All four operators are
syntactic. Semantic transformations exploiting dependency schemes
or quantifier-alternation depth could potentially activate deeper
solver code paths.


\section{Future Work}\label{sec:future}

Several directions follow naturally from this work.
\emph{Larger benchmark sets} from QBFLIB or the QBF Gallery would
enable statistically meaningful per-family comparisons.
\emph{Additional operators} could target clause learning
(Q-resolution), dependency scheme computation, or certificate
generation. \emph{Mutation testing} of the solver source itself
would quantify the sensitivity of the test method by checking
whether injected bugs are detected by the equivalence flag.
\emph{Automated operator synthesis} using constraint solving or
program synthesis could discover new truth-value-preserving
transformations beyond those designed by hand. Finally,
\emph{cross-solver coverage} via Rust-native tooling
(e.g.\ \texttt{cargo-llvm-cov}) would allow the DepQBF/CAQE runtime
asymmetry to be investigated symmetrically --- arguably the most
important open question left by this thesis.